\documentclass[11pt,a4paper]{article}
\usepackage[margin=25mm]{geometry}
\usepackage{amsmath,amssymb,amsthm,hyperref}
\setlength{\emergencystretch}{2em}
\newtheorem{theorem}{Theorem}
\newcommand{\gen}[1]{\langle#1\rangle}
\title{One deficit at an odd singleton end}
\author{Human Compression project: L044 working continuation}
\date{}
\begin{document}\maketitle
This is a refinement of the odd-degree equality interface in source
Proposition 5.31, not a recovery of its exceptional diagram.
Let $G=\gen{\rho_0,\ldots,\rho_{r-1}}\le A_n$ be a transitive
connected string C-group, $n=2r+1\ge13$, with all maximal omissions
intransitive. Suppose $G_0$ has orbits $\{a\}$ and its complement.
Put $b=a^{\rho_0}$, $H=G_{\ge2}$, $K=G_{>2}$,
$\sigma=\rho_1$, and $\rho_0=(a\ b)\delta$.
Thus $H$ fixes $a,b$ and centralizes the odd involution $\delta$
on $X=\Omega\setminus\{a,b\}$.
Suppose every colour $j\ge2$ has at least two ambient-eligible edges.
The selected-localization lemmas and the full retained odd-degree
maximal-string theorem are required below.

\begin{theorem}
Either $H$ fixes at least three points, so that
\[
 1\le \frac{n-|\operatorname{supp}(G_{>1})|-1}{2},       \tag{1}
\]
or its only nontrivial orbits are a two-point set $E=\{u,v\}$
and an odd set $V$ of size $n-4$, and
\[
               \rho_0=(a\ b)(u\ v).                    \tag{2}
\]
No claim that (2) already determines the complete source diagram is made.
\end{theorem}
\begin{proof}
Let $f$ count all fixed points of $H$ and $c$ its nontrivial orbits.
The ambient interval bound gives $f+c\le5$. If $f\ge3$, (1)
is immediate. Suppose $f=2$. Then $c\le3$, and $X$ is the
union of those nontrivial orbits, of odd total size.
A commuting involution on a preserved transitive orbit is either
free or the identity. If $c=1$, odd size forces $\delta=1$,
contradicting oddness. If $c=2$, the orbits cannot be exchanged,
again by total parity. Exactly one is even and carries the odd free
action of $\delta$; write its size as $2q$ with $q$ odd.
The other orbit $V$ is odd and fixed by $\delta$.
The case $q=1$ is precisely (2). We exclude $q\ge3$ below.

First dispose of $c=3$. If two orbits are exchanged, their common
size $m$ is odd and the third orbit $V$ is odd and fixed. The rank
identity is $r-2=m+(|V|-3)/2$. The exchanged-row argument from the
exceptional global-pair proof applies verbatim after replacing
$\rho_3,\rho_4,5$ by $\rho_0,\rho_1,2$: row-eligible labels have
at most $m-1$ members; every other label is confined to $V$ and
has at most $(|V|-1)/2$ members. Equality gives the ordinary
consecutive $S_m$ row path and a supported subgroup transitive
on $V$, using the no-one-colour-hole bound.
If the first tail label belongs to the row path, the latter subgroup
centralizes $\sigma$, forcing $V$ invariant under $G_0$.
Otherwise the full row path belongs to $K$, so $K$ fixes no point
of either row. Since $x=b^\sigma\ne b$ is $K$-fixed and
$\sigma$ fixes $a$, it lies in $V$. This supplies the one fact
for which the exceptional-head proof used its pure-head lemma.
Its terminal-component and last-row-label argument then applies
unchanged and contradicts transitivity of $G_0$.

If all three orbits are preserved, an odd free action occurs on
at least one. With exactly one free orbit $O$ of size $2q$, split
labels into those with any eligible edge in $O$ and those with none.
The first set has at most $q-1$ labels by the projected pair forest.
The second is confined to the two fixed orbits of total odd size $v$,
so the selected-localization bound gives at most $(v-3)/2$ labels.
Their sum is strictly below the required $r-2=q+(v-3)/2$.
If there are two free orbits of sizes $2q_1,2q_2$ and one fixed
odd orbit $V$ of size $v$, assign a label first to the first free
orbit if possible, next to the second if possible, and otherwise
to $V$. This is a disjoint partition. The pair forests and localized
bound give respectively $q_1-1,q_2-1,(v-1)/2$ labels, whose sum
is strictly below $r-2=q_1+q_2+(v-3)/2$.
Three free orbits have even total size and are impossible.
This excludes $c=3$.

Return to $c=2$ with a free orbit $O$ of size $2q$, $q\ge3$
odd, and a fixed odd orbit $V$ of size $v$. Let $L$ be all tail
labels with an eligible edge in $O$, and $F$ all remaining labels.
The pair forest gives $|L|\le q-1$. All eligible $F$-edges lie in
$V$, so selected localization gives $|F|\le(v-1)/2$.
Since $r-2=q+(v-3)/2$, equality holds in both.
The $q-1$ pair images of $L$ are independent and form the ordinary
$S_q$ path. Independence follows because an orbit path among pairs
lifts using $\delta\in G_j$ and contradicts the chosen ambient cut.
The consecutive path cuts then force every $F$ label to have trivial
pair image. Thus each acts on $O$ only by individual pair swaps,
whose edges are ambient-ineligible.

The set $L$ is consecutive with $|L|\ge2$, so $F$ has no
one-colour hole. All its eligible edges are confined to $V$;
the corresponding orbit bound and $2|F|=v-1$ force $G_F$
transitive on $V$. Its orbits in $O$ have size at most two.
If $2\in L$, then $G_F\le K$ centralizes $\sigma$; its unique
orbit of size $v\ge3$ is $V$, and $\sigma$ preserves it.
This contradicts transitivity of $G_0$.

Suppose $2\in F$. The whole $L$ path belongs to $K$, which is
transitive on the $q$ pair-blocks in $O$, and hence fixes no point
there. The $K$-fixed point $x=b^\sigma\ne b$ must be in $V$.
Any terminal component of $F$ is contained in $K$, fixes $x$,
and commutes with its initial component. As $G_F$ is transitive
on $V$, that terminal component fixes $V$ pointwise. Its only
remaining edges are individual $\delta$-pair edges in $O$.
They are all ineligible, contradicting the hypothesis on each
of its labels. Therefore the terminal component is empty.

The last label $\lambda$ of the now terminal $L$ path centralizes
$G_F$ and fixes $x$, so fixes $V$ pointwise. Each $K$-orbit in
$O$ meets every pair-block, and $\lambda$ has a nontrivial pair
transposition, so acts nontrivially on each such orbit. It fixes
all points outside $O$. Since $\sigma$ centralizes $K$, it must
preserve $O$, again contradicting transitivity of $G_0$.
This eliminates $q\ge3$ and proves the theorem.
\end{proof}

\begin{theorem}[An actual fixed-support budget]
Under the same hypotheses, either $H=G_{>1}$ fixes at least three
points or $K=G_{>2}$ fixes at least five. Consequently there is
$h\in\{1,2\}$ such that, with $X_h=\operatorname{supp}(G_{>h})$,
\[
                         2h\le n-|X_h|-1.              \tag{3}
\]
\end{theorem}
\begin{proof}
Only residual (2) needs consideration. Write $x=b^\sigma$.
Transitivity of $G_0=\gen{\sigma,H}$ gives $x\ne b$, and
$\sigma$ fixes $a$, so $x\in E\cup V$. The group $K$ preserves
$E$ and $V$, fixes $a,b,x$, and preserves $E^\sigma$ because
it centralizes $\sigma$.

Suppose first that $K$ fixes $E$ pointwise. It then fixes
$E^\sigma$ pointwise too. If $x\in V$, the five distinct points
$a,b,u,v,x$ are fixed. If $x\in E$, rename it $u$.
Then $u^\sigma=b$. The other image $w=v^\sigma$ must lie in
$V$: it cannot be $a,b,u$, and if it were $v$, then both $H$
and $\sigma$ would preserve $\{b,u,v\}$, contrary to transitivity
of $G_0$ on its larger complementary orbit. Thus the five points
$a,b,u,v,w$ are fixed by $K$.

Suppose now that $K$ does not fix $E$ pointwise. Then $x\notin E$,
for its being $K$-fixed would force both points fixed. Also
$E^\sigma\ne E$, since otherwise $\gen{\sigma,H}$ would
preserve $E$. Their intersection cannot have size one: it would
be a $K$-fixed singleton, again forcing $K$ to fix $E$.
They are therefore disjoint. Since $x\in V$, neither point of
$E^\sigma$ is $b$; neither is $x$ because $x^\sigma=b$.
Write $E^\sigma=\{w,z\}\subset V\setminus\{x\}$ and put
$Y=V\setminus\{x,w,z\}$. This is $K$-invariant.

For every colour $j>2$, an edge outside $Y$ is either an $E$-edge
or an $E^\sigma$-edge; the remaining points $a,b,x$ are fixed.
The first kind is ambient-ineligible since $\rho_0\in G_j$.
The second kind is likewise ineligible since
$\rho_0^\sigma=(a\ x)(w\ z)\in G_j$.
Thus all eligible edges of the nonempty full interval $\{3,\ldots,r-1\}$
are confined to $Y$, and each colour has two by hypothesis.
But $|V|=n-4$ and $n=2r+1$ give
\[
               |Y|=n-7=2(r-3).
\]
The localized interval bound gives $2(r-3)\le |Y|-o_K(Y)<|Y|$,
a contradiction. This proves the five-point fixed set and hence (3).
All fixed singleton orbits of $K$ in $Y$ are retained in this count.
\end{proof}

\paragraph{Attribution and exact unfinished boundary.}
The source is Cameron--Fernandes--Leemans--Mixer,
\emph{Highest rank of a polytope for $A_n$}, arXiv:1605.09173v1,
Proposition 5.31, physical pages 42--44. Its odd-residual and
fixed-support strategy is reused. The full accepted exceptional
 global-pair proof, the two selected-localization theorems and all
maximal-string inputs remain paid. The new argument adapts these to
$n=2r+1$ and resolves the literal two-point residual by actual
ambient confinement, yielding (3).
The proof does not identify the complete source equality graph or
assert its stronger alternatives at other ranks. Nor does it justify
inheriting ambient eligibility in an arbitrary parabolic: that must
be checked in the proposed embedding. No new computation, projected
IP, formalization or original credit is involved.
\end{document}
